\section{电磁对偶与单极子}\label{sec:emdual}\index{ElectromagneticDuality@电磁对偶|hyperpage}
麦克斯韦理论中交换电场与磁场的电磁对偶，在本讲义中发挥着核心作用。因而，我们不妨先在不引入超对称性这一额外复杂因素的情况下，对它作一回顾。我们也将回顾磁单极子的基本性质。
\subsection{电荷与磁荷}\label{charges}\index{ElectricCharge@电荷|hyperpage}\index{MagneticCharge@磁荷|hyperpage}
考虑一个 $\U(1)$ 规范场，它由规范势 $A=A_\mu dx^\mu$ 和场强 $F=\frac12F_{\mu\nu} dx^\mu \wedge dx^\nu$ 描述，其中 $F_{\mu\nu}=\partial_{[\mu} A_{\nu]}$。场强在规范变换
\begin{equation}
A\to A+i g^{-1} dg
\end{equation}
下不变，其中 $g$ 是从时空到绝对值为一的复数的映射，即 $|g|=1$。我们可以用实函数 $\chi$ 将其写为 $g=e^{i\chi}$，于是得到更熟悉的形式
\begin{equation}
A\to A-d\chi,\label{agauge}
\end{equation}
不过，对我们而言，重要的一点是 $\chi$ 可以是多值的，因此我们作如下等同：
\begin{equation}
\chi\sim \chi+2\pi. \label{chiperiod}
\end{equation}

考虑一个场 $\phi$，其规范变换为
\begin{equation}
\phi \to g^n \phi
\end{equation}
这里，我们要求 $g$ 唯一地确定系统中所有场的变换。因此，$n$ 必须是整数；分数次幂的定义并不唯一。

协变导数为
\begin{equation}
D_\mu\phi = \partial_\mu\phi + inA_\mu \phi,
\end{equation}
而动能项 $|D_\mu\phi|^2$ 是规范不变的。我们将规范场的作用量写为
\begin{equation}
S_\text{Maxwell}=\int d^4 x \frac1{2e^2} F_{\mu\nu} F_{\mu\nu}.\label{maxwell}
\end{equation}
分母中的系数 $2$ 稍显不合惯例，但这一选择会消去后面出现的各种烦人的因子。于是，对场 $\phi$ 量子化得到的两个粒子之间的力，正比于 $e^2n^2$。在唯象学文献中，组合 $en$ 经常被称为电荷，但在本讲义中，我们把整数 $n$ 称为电荷。人们也可能倾向于对 $F$ 作重新标度，以消去上式分母中的因子 $e^2$。不过，我们坚持采用 $\chi$ 的周期为 $2\pi$ 的约定，参见式\eqref{chiperiod}。

在一次量子化框架中，经过 Wick 转动到欧氏号差之后，一个电荷为 $n$ 的带电粒子通过
\begin{equation}
S= in \int_L dx^\mu A_\mu\label{particle}
\end{equation}
与规范场耦合，其中 $L$ 是世界线。在这一描述中，$n$ 的整数性可以如下看出。由于 $\chi$ 具有式\eqref{chiperiod}所示的周期性，线积分 $\int dx^\mu A_\mu$ 只能确定到相差 $2\pi$ 的整数倍。在路径积分中，$e^{iS}$ 必须有良好定义。因而，$n$ 必须是整数。

将式\eqref{maxwell}与式\eqref{particle}相加，并写出 $A_\mu$ 的运动方程，可得
\begin{equation}
\int_{S^2}  \frac{4\pi}{e^2}\vec E \cdot d\vec n=
\int_{S^2}  \frac{4\pi}{e^2} \star F=2\pi n,\label{equantize}
\end{equation}
其中 $E_i=F_{0i}$ 是电场的分量，$S^2$ 是无穷远处的球面，而
\begin{equation}
\star F=\frac12 (\star F)_{\mu\nu} dx^\mu\wedge dx^\nu
\end{equation}
中的
\begin{equation}
(\star F)_{\mu\nu}=\frac12\epsilon_{\mu\nu\rho\sigma} F^{\rho\sigma}
\end{equation}
为对偶场强。我们也会交替使用记号 $\tilde F=\star F$。

接下来，考虑去掉原点的空间。用一个球面包围原点。北、南半球上的规范场 $A_{N,S}$ 在赤道上通过规范变换联系起来：
\begin{equation}
A_N = A_S+ig^{-1} dg
\end{equation}
于是有
\begin{equation}
\int_{S^2} F = \int_{N} F + \int_{S} F = \int_\text{赤道}(A_N-A_S) =\int_{\theta=0}^{\theta =2\pi} \frac{d\chi}{d\theta} d\theta=2\pi m,\label{mquantize}
\end{equation}
其中 $m$ 为整数。我们把 $m$ 称为这一组态的磁荷。库仑型磁场中所含的能量在原点处发散；不过，不必对此过分担心，因为量子化的带电粒子也具有库仑型电场，而该电场的能量同样发散。通过重整化，两者都会变为有限。当 $m$ 非零时，这一组态称为磁单极子。通常，我们简称它为单极子。

\begin{figure}[h]
\[
\inc{angular}
\]
\caption{带电粒子与带磁荷粒子同时存在时产生的角动量。实线箭头、虚线箭头和双线箭头分别表示电场、磁场和坡印廷矢量。\label{fig:angular}}
\end{figure}

在两个分离的点上，分别放置一个电荷为 $n$ 的粒子和一个磁荷为 $m$ 的粒子。两者的电场和磁场共同通过坡印廷矢量，产生绕两点连线方向的角动量，参见图\ref{fig:angular}。仔细计算表明，电磁场中所含的总角动量为 $\hbar nm/2$，这与量子力学中的角动量量子化相一致。

更一般地，我们可以考虑双荷子，即同时具有电荷和磁荷的粒子。如果有一个粒子的电荷为 $n$、磁荷为 $m$，另一个粒子的电荷为 $n'$、磁荷为 $m'$，那么总角动量就是 $\hbar/2$ 乘以
\begin{equation}
 nm'-mn'.
\end{equation}
我们把这一组合称为两组荷 $(n,m)$ 和 $(n',m')$ 的狄拉克配对\index{DiracPairing@狄拉克配对|hyperpage}。

\subsection{S 变换与 T 变换}\label{ST}\index{STransformation@S 变换|hyperpage}\index{TTransformation@T 变换|hyperpage}\index{ComplexifiedCoupling@复化耦合|hyperpage}
麦克斯韦方程为
\begin{equation}
\partial_{[\mu}F_{\nu\rho]}=0,\qquad
\partial_{\mu}F_{\mu\nu}=0
\end{equation}
或等价地，用微分形式记号写为
\begin{equation}
dF=d\star F=0.
\end{equation}
这组方程在交换
\begin{equation}
F\leftrightarrow \star F.
\end{equation}
下不变。用电场 $\vec E$ 和磁场 $\vec B$ 表示，并示意性地记为 $F=(\vec E,\vec B)$，则这一变换的作用为
\begin{equation}
F=(\vec E,\vec B) \longrightarrow
\star F=(\vec B,-\vec E) \longrightarrow
\star^2 F=(-\vec E,-\vec B).
\end{equation}
这一操作经常被称为 $S$ 变换。

为保持式\eqref{equantize}、\eqref{mquantize}中的电荷与磁荷量子化，对偶场强 $F_D$ 和对偶耦合 $e_D$ 必须定义为满足
\begin{equation}
F_D=\frac{4\pi}{e^2}{}\star F,\qquad
\frac{4\pi}{e^2}\frac{4\pi}{e_D^2}=1.
\end{equation}

在这一变换下，荷 $(n,m)$ 按照如下方式变换：
\begin{equation}
\begin{array}{r|ccc}
\text{粒子 1} & (n,m) & \stackrel{S}{\longrightarrow} & (-m,n), \\
\text{粒子 2} &(n',m') & \stackrel{S}{\longrightarrow} & (-m',n'), \\
\text{狄拉克配对} & nm'-n'm & = &  -mn'-(-m')n.
\end{array}
\end{equation}
注意，狄拉克配对在这一操作下保持不变。

假设存在一个中性的实标量场 $\phi$，而 $\U(1)$ 规范场的作用量由下式给出：
\begin{equation}
\frac{1}{2e(\phi)^2} F_{\mu\nu}F_{\mu\nu} + \frac{\theta(\phi)}{16\pi^2} F_{\mu\nu} \tilde F_{\mu\nu}.
\label{u1coupling}
\end{equation}
现在，麦克斯韦方程为
\begin{align}
\partial_{[\mu}F_{\nu\rho]}&=0,\label{MaxwellEq1}\\
\partial_{\mu}\left[ \frac{4\pi}{e(\phi)^2} F_{\mu\nu}+\frac{\theta(\phi)}{2\pi}\tilde F_{\mu\nu}\right]&=0.
\label{MaxwellEq2}
\end{align}
像前面一样作分解 $F=(\vec E,\vec B)$。上述方程表明，满足高斯定律的磁场仍为 $\vec B$，但满足高斯定律的电场现在是如下组合：
\begin{equation}
\vec E_\text{conserved} = \frac{4\pi }{e(\phi)^2} \vec E+\frac{\theta(\phi)}{2\pi} \vec B.
\end{equation}
因此有
\begin{equation}
\int_{S^2} \vec B \cdot d\vec n= 2\pi m,\quad
\int_{S^2} \vec E_\text{conserved} \cdot d\vec n = 2\pi n
\end{equation}
其中 $m$ 和 $n$ 是第\ref{charges}小节引入的整数。这揭示了一个有趣的事实：让我们绝热地改变 $\phi$，从而改变 $\theta(\phi)$。由于 $n$ 是整数，它不能改变。因而，为保持 $\vec E_\text{conserved}$ 不变，$\vec E$ 会获得一个正比于 $\theta(\phi)\vec B$ 的贡献。这称为 Witten 效应\index{WittenEffect@Witten 效应|hyperpage}\cite{Witten:1979ey}。

于是，S 变换交换的是 $(\vec E)_\text{conserved}$ 与 $\vec B$。对偶规范场强 $F_D$ 因此应为
\begin{equation}
F_D=\frac{4\pi}{e(\phi)^2} {}\star F - \frac{\theta}{2\pi}F
\end{equation}
使得式\eqref{MaxwellEq2}恰好成为
\begin{equation}
\partial_{[\mu}F_{D,\nu\rho]}=0.
\end{equation}
我们发现，式\eqref{MaxwellEq1}变为
\begin{equation}
\partial_{\mu}\left[ \frac{4\pi}{e_D(\phi)^2} F_{D,\mu\nu}+\frac{\theta_D(\phi)}{2\pi}\tilde F_{D,\mu\nu}\right]=0
\end{equation}
其中 $e_D(\phi)$、$\theta_D(\phi)$ 由
\begin{equation}
\tau_D(\phi)  = -\frac{1}{\tau(\phi)}
\end{equation}
确定，而
\begin{equation}
\tau(\phi)= \frac{4\pi i}{e(\phi)^2} + \frac{\theta(\phi)}{2\pi},\qquad
\tau_D(\phi)= \frac{4\pi i}{e_D(\phi)^2} + \frac{\theta_D(\phi)}{2\pi}.\label{u1tau}
\end{equation}
这一组合 $\tau(\phi)$ 称为复化耦合。

我们还知道，在量子理论中，$\theta(\phi)$ 与 $\theta(\phi)+2\pi$ 无法区分，因为欧氏路径积分中的被积函数仅改变一个因子
\begin{equation}
\exp\left[i\int d^4x \frac{1}{8\pi}  F_{\mu\nu} \tilde F_{\mu\nu}\right]
\end{equation}
而该因子始终为一\footnote{严格地说，这仅在自旋流形上成立。注意，$\int d^4x (8\pi)^{-1}  F_{\mu\nu} \tilde F_{\mu\nu}= \pi \int (F/2\pi)^2=\pi \int c_1(F)^2$。在自旋流形上，交叉形式为偶形式，最后一个表达式是 $2\pi$ 的整数倍。有关非自旋流形上的细微问题，参见文献\cite{Witten:1995gf}。}。我们将其称为 $T$ 变换。不过，它确实会通过加上 $\vec B$ 而改变 $\vec E_\text{conserved}$。等价地，它使荷 $(n,m)$ 发生如下变化：
\begin{equation}
\begin{array}{r|ccc}
\text{粒子 1} & (n,m) & \stackrel{T}{\longrightarrow} & (n+m,m), \\
\text{粒子 2} &(n',m') & \stackrel{T}{\longrightarrow} & (n'+m',m'), \\
\text{狄拉克配对} & nm'-n'm & = &  (n+m)m'-(n'+m')m.
\end{array}
\end{equation}
可以看到，两个粒子的狄拉克配对保持不变。它对复化耦合 $\tau(\phi)$ 的作用为
\begin{equation}
\tau(\phi)_\text{old}  \stackrel{T}{\longrightarrow} \tau(\phi)_\text{new} = \tau(\phi)_\text{old}+1.
\end{equation}
变换 $S$ 与 $T$ 生成 $\SL(2,\bZ)$ 对荷 $(n,m)$ 的作用：
\begin{align}
S&=\begin{pmatrix}
0&-1\\1&0
\end{pmatrix}, &
S \begin{pmatrix}n \\ &=
\begin{pmatrix}-m \\,&
S\tau &= -\frac1\tau \label{S}\\
T&=\begin{pmatrix}1&1\\0&1\end{pmatrix}, &
T \begin{pmatrix}n \\ &=
\begin{pmatrix}n+m \\,&
T\tau &= \tau+1. \label{T}
\end{align}
一般而言，它对 $\tau$ 的作用为分式线性变换
\begin{equation}
\begin{pmatrix}a&b \\c&d\end{pmatrix} \in \SL(2,\bZ) :
\tau \to \frac{d\tau+b}{c\tau+a}.
\end{equation}

\subsection{'t Hooft-Polyakov 单极子}\label{sec:monopole}\index{HooftPolyakovMonopole@'t Hooft-Polyakov 单极子|hyperpage}
这里，我们总结磁单极子的一些性质，后面将反复引用这些性质。有关本小节所讨论内容的详细阐述，读者可参阅文献\cite{Harvey:1996ur,Weinberg:2006rq}等综述，或教科书\cite{Shifman}。Coleman 的综述\cite{Coleman:1982cx}也非常具有启发性。\footnote{遗憾的是，这篇综述并未收入《Aspects of Symmetry》一书。R.~Stora 所译的法文版本也可见于文献\cite{Coleman1982cxF}，其排版比文献\cite{Coleman:1982cx}中的版本精美得多。}
\subsubsection{经典性质}\index{BPSBound@BPS 界|hyperpage}\index{BPSEquation@BPS 方程|hyperpage}
考虑一个含有伴随表示标量场的 $\SU(2)$ 规范理论，其作用量为
\begin{equation}
S=\int d^4 x\frac{1}{g^2}\left[\frac{1}{2}\tr F_{\mu\nu} F_{\mu\nu} + \tr D_\mu\Phi D_\mu\Phi \right].
\end{equation}
场 $\Phi$ 是一个无迹厄米 $2\times 2$ 矩阵。

考虑满足
\begin{equation}
\Phi= \begin{pmatrix}a & 0 \\0 & -a\end{pmatrix}.\label{phivev}
\end{equation}
的真空。当 $a\neq 0$ 时，$\SU(2)$ 规范对称性破缺为 $\U(1)$。事实上，真空期望值\eqref{phivev}与具有如下形式的规范场强对易：
\begin{equation}
F_{\mu\nu}=\begin{pmatrix}
F_{\mu\nu}^{\U(1)} & 0\\
0 & -F_{\mu\nu}^{\U(1)}
\end{pmatrix}\label{u1embed}
\end{equation}
其中 $F_{\mu\nu}^{\U(1)}$ 是按照第\ref{charges}小节约定归一化的 $\U(1)$ 规范场强。注意，标量场 $\Phi$ 的非对角分量的量子，相对于这一 $\U(1)$ 场具有电荷 $\pm 2$；这可以通过展开协变导数得出。

我们考虑的是规范理论，因此场 $\Phi$ 不必严格取为式\eqref{phivev}右端的形式。相反，我们只需要 $\Phi$ 的本征值为 $\pm a$。因而，可以考虑如下形式的组态：
\begin{equation}
\Phi(x)=\frac{x_i \sigma^i}{|x|} f(|x|) a
\label{tHP}
\end{equation}
其中 $i=1,2,3$，而 $f(r)$ 是满足
\begin{equation}
\lim_{r\to 0} f(r)=0,\qquad
\lim_{r\to \infty} f(r)=1.
\end{equation}
的无量纲函数。在空间无穷远处，$\Phi$ 的真空期望值与式\eqref{phivev}共轭，因此这一组态可被视为式\eqref{phivev}所给真空的激发。

$\SU(2)$ 中未破缺的 $\U(1)$ 沿着 $\Phi$ 的方向。因而，至少在 $r\gg 0$ 时，$\U(1)$ 场强 $F_{\mu\nu}^{\U(1)}$ 的一个更一般的定义是如下组合：
\begin{equation}
F_{\mu\nu}^{\U(1)}:=\frac{1}{2a}\tr F_{\mu\nu} \Phi .\label{gauge-inv-def-of-u1}
\end{equation}
在 $r\gg 0$ 的区域中，尝试通过一个规范变换将组态\eqref{tHP}化为式\eqref{phivev}。除南极点以外，这可以光滑地实现，所用的规范变换为
\begin{equation}
\exp[ i\frac{\varphi}{2} (-\sigma^1 \sin\theta + \sigma^2 \cos\theta )],\quad
\text{其中}\quad \frac{\vec x}{|x|}=(\cos\varphi,\sin\varphi\cos\theta,\sin\varphi\sin\theta).
\end{equation}
由此得到的南极点附近的规范变换为
\begin{equation}
i (-\sigma^1 \sin\theta + \sigma^2\cos\theta) =  \exp[ -i\theta \sigma^3 ] \cdot (i\sigma^2).
\end{equation}
当 $\theta$ 从 $0$ 变到 $2\pi$ 时，可以看到 $\U(1)$ 场 $F^{\U(1)}_{\mu\nu}$ 的磁荷为 $m=1$，因此它是一个单极子。这一结果最初由 't Hooft 和 Polyakov 发现。注意，它与场 $\Phi$ 的粒子的狄拉克配对为 2，是最小允许值的两倍。

我们来计算这一场组态中所含的能量。动能为 $1/g^2$ 乘以下式：
\begin{align}
\int d^3x \left[\tr B_{i} B_i + \tr D_i\Phi D_i\Phi \right]
&=\int d^3 x \left[ \tr(B_i\mp D_i\Phi)^2 \pm 2 \tr B_i D_i\Phi \right]\\
&\ge \pm 2 \int d^3 x  \tr B_i D_i\Phi = \pm 2 \int d^3 x \partial_i \tr B_i\Phi \\
&= \pm 2 \int_{S^2}  d\vec n \cdot \tr \vec B\Phi
\end{align}
其中最后一个积分是在空间无穷远处的球面上进行的。根据式\eqref{gauge-inv-def-of-u1}，它等于 $\pm 2(2a)(2\pi m)$，其中 $m$ 为磁荷。因而得到下界
\begin{equation}
\text{（单极子的能量）}\ge \frac{4\pi }{g^2} (2a) |m|  \label{monopole-energy}
\end{equation}
这称为 Bogomolnyi-Prasad-Sommerfield（BPS）界。当且仅当
\begin{equation}
B_i=\pm D_i\Phi,
\end{equation}
时，该不等式取等号；这一方程称为 BPS 方程。它确定了式\eqref{tHP}中函数 $f(r)$ 的形式。

\subsubsection{半经典性质}\index{ZeroModes@零模|hyperpage}\index{Dyon@双荷子|hyperpage}\index{SemiclassicalQuantization@半经典量子化|hyperpage}\index{FlavorSpinor@味旋量|hyperpage}
给定这样一个显式的单极子解，我们可以利用对称性构造与它相关的其他解。首先，组态\eqref{tHP}的中心位于坐标系原点。可以把单极子的中心平移到空间 $\bR^3$ 中的任意一点 $\vec y$。由此得到三个零模。

另一个零模由如下规范变换得到：
\begin{equation}
e^{ i \alpha \Phi/a }\label{gauge-zm}
\end{equation}
注意，在无穷远处为零的规范变换是物理系统的一种冗余，而在无穷远处不为零的规范变换则被认为会改变经典组态。对一般的 $\alpha$，变换\eqref{gauge-zm}会改变 $F_{ij}(x)$ 的渐近行为；但在 $\alpha=\pi$ 时，这一变换对伴随表示中的场作用平凡。因此，$\alpha$ 是一个取值范围为 $0\le\alpha<\pi$ 的角变量。

单极子的半经典量子化涉及非零模的 Fock 空间，以及依赖于零模 $\vec y$ 和 $\alpha$ 的波函数 $\psi(\vec y,\alpha)$。沿 $\vec y$ 方向的波函数描述单极子质心的空间运动。沿 $\alpha$ 方向的波函数则描述单极子的电荷，这一点可以如下看出。

将式\eqref{gauge-zm}与式\eqref{u1embed}、\eqref{gauge-inv-def-of-u1}比较，可以看到，未破缺的 $\U(1)$ 整体规范变换 $e^{i\varphi}$ 使 $\alpha$ 平移为
\begin{equation}
\alpha\to \alpha+\varphi.\label{shift-el}
\end{equation}
回顾一下，电荷为 $n$ 的态 $\ket\psi$ 在 $\U(1)$ 整体变换 $e^{i\varphi}$ 下按照
\begin{equation}
\ket\psi \to e^{in\varphi} \ket\psi.\label{www}
\end{equation}
变换。现在，由于 $\alpha$ 是周期为 $\pi$ 的变量，$\psi(\alpha)$ 可以展开为 $e^{i2d\alpha}$ 的线性组合，其中 $d$ 为整数。在变换\eqref{shift-el}下，波函数按照式\eqref{www}中 $n=2d$ 的情形变换，因此，具有这一零模波函数的单极子态的电荷为 $2d$。

总之，由半经典量子化得到的电荷与磁荷组合 $(n,m)$ 具有 $(n,m)=(2d,1)$ 的形式，其中 $d$ 为整数。这一结果最初由 Julia 和 Zee 发现：一旦对 't Hooft-Polyakov 单极子进行量子化，我们得到的就不仅是一个纯磁单极子，而是从 $d=-\infty$ 到 $+\infty$ 的一整塔双荷子态。

最后，考虑 't Hooft-Polyakov 单极子\eqref{tHP}中的费米子零模的效应。首先考虑伴随表示中的两个 Weyl 费米子 $\lambda$、$\tilde\lambda$，其拉格朗日量为
\begin{equation}
\tr \bar\lambda \slash D \lambda + \tr\bar{\tilde\lambda} \slash D \tilde\lambda
+ c(\tr  \lambda [\Phi,  \tilde\lambda ] + \tr \bar \lambda [\Phi,\bar{\tilde\lambda}]).
\end{equation}
将协变导数 $D$ 中的规范势以及标量场都视为背景，并把 $\lambda$、$\tilde\lambda$ 分解成角动量本征态。轨道角动量的下界由狄拉克配对给出，在这里等于 $\hbar$。旋量场的自旋为 $\hbar/2$。因此，最低角动量的态具有 $\hbar/2$ 的自旋。已知当系数 $c$ 取值于某一区间时，存在一对零模 $b_\alpha$，其中 $\alpha=1,2$ 是空间转动群 $\SO(3)$ 的旋量指标。半经典量子化将它们提升为一对费米子振子：
\begin{equation}
\{b_\alpha,b_\beta^\dagger\}=\delta_{\alpha\beta}.
\end{equation}
从玻色部分半经典量子化得到的一个态 $\ket\psi$ 出发，它们产生四个态：
\begin{equation}
\begin{array}{ccccc}
&\raisebox{-.8ex}{\rotatebox{30}{$\leftrightarrow$}}& b^\dagger_1 \ket\psi & \leftrightarrow & b^\dagger_1 b^\dagger_2 \ket\psi \\
\ket\psi & \leftrightarrow & b^\dagger_2\ket\psi & \rotatebox{30}{$\leftrightarrow$}
\end{array}.
\end{equation}
这相当于一个复玻色子和一个费米子。

假设再引入一对伴随 Weyl 费米子 $\lambda'$、$\tilde\lambda'$。那么将得到另一对费米子振子 $b_\alpha'$。它们共同生成 $2^4=16$ 个态，包括一个有质量矢量（有 3 个态）、四个有质量旋量（有 8 个态）以及五个有质量标量。

接下来，考虑二重态表示中的 $2N$ 个 Weyl 费米子 $\psi_i^a$，其中 $a=1,2$，$i=1,\ldots,2N$，其拉格朗日量为
\begin{equation}
 \bar\psi_i \slash D \psi_i
+ c'  (\psi^a_i \Phi_{(ab)}  \psi^b_i +\bar\psi^a_i \Phi_{(ab)}  \bar\psi^b_i ).
\end{equation}
注意，该拉格朗日量具有作用于指标 $i$ 的 $\SO(2N)$ 味对称性。

现在，$\psi$、$\tilde\psi$ 的量子相对于未破缺的 $\U(1)$ 的电荷为 $1$。因而，狄拉克配对为 $\hbar/2$。与内禀自旋 $\hbar/2$ 作张量积，可知最小轨道角动量为 $0$。已知适当选择 $c'$ 后，该费米子系统具有零模 $\gamma_i$，$i=1,\ldots,2N$。半经典量子化之后，它们成为一组满足如下对易关系的费米子算符：
\begin{equation}
\{\gamma_i,\gamma_j\}=\delta_{ij}.
\end{equation}
这就是 $\SO(2N)$ 的伽马矩阵的对易关系。单极子态构成 $\gamma_i$ 的表示，这意味着它们按照味对称群 $\SO(2N)$ 的旋量表示变换。

$\SU(2)$ 规范对称性二重态表示中的场还有另一个效应。即在规范零模\eqref{gauge-zm}中，$\alpha=\pi$ 给出矩阵
\begin{equation}
\begin{pmatrix}-1 & 0 \\0 & -1\end{pmatrix}
\end{equation}
它对二重态表示中的场具有非平凡作用。于是，规范零模 $\alpha$ 的周期现在是 $2\pi$，沿 $\alpha$ 方向的波函数可以取为任意整数 $n$ 对应的 $e^{in\alpha}$。因此，电荷 $n$ 既可以为偶数，也可以为奇数。算符 $\gamma_i$ 来自二重态表示中场的模，因此会使电荷改变 $\pm1$。

可以通过
\begin{equation}
\Gamma=\gamma_1\gamma_2 \cdots \gamma_{2N},
\end{equation}
定义味旋量的手征性 $\Gamma$，从而将 $\SO(2N)$ 的旋量分解为正手征和负手征旋量。算符 $\gamma_i$ 的作用会改变味旋量的手征性。结合上一段所见的 $\U(1)$ 电荷的行为，我们得出结论：单极子态的 $\U(1)$ 电荷的奇偶性，与味旋量表示的手征性相关。
